\phantomsection
\section*{ВВЕДЕНИЕ}
\addcontentsline{toc}{section}{Введение}

Опоры скольжения остаются базовым элементом подавляющего большинства роторных машин: турбоагрегатов, нагнетательных и буровых насосов, центробежных и винтовых компрессоров нефтегазового профиля. Именно через эти узлы передаётся радиальная нагрузка на вращающийся вал, и от стабильности их работы зависит возможность сохранения проектных режимов всего агрегата. Несущий жидкостный слой, разделяющий цапфу и вкладыш, образуется только при определённом сочетании частоты вращения, вязкости и зазора~-- именно при этих условиях возникает гидродинамическое давление, исключающее металлический контакт поверхностей.

Отказ подшипника на оборудовании добычи и транспорта углеводородов сопровождается не только остановом дорогостоящего агрегата, но и потенциальным экологическим ущербом, особенно на удалённых и шельфовых объектах. По этой причине задачи, направленные на повышение нагрузочной способности и снижение энергетических потерь во вкладышах скольжения, относятся к числу приоритетных в трибологии нефтегазового машиностроения.

Один из наиболее активно развиваемых сегодня подходов~-- модификация рабочей поверхности подшипника регулярным микрорельефом, формируемым лазерной обработкой. Современная лазерная техника позволяет получать массивы микроуглублений строго заданной формы, размеров и плотности расположения, что открывает возможность управления трибологическими характеристиками непосредственно на этапе изготовления. Опубликованные экспериментальные данные демонстрируют снижение коэффициента трения на 10-30\,\% и одновременный прирост несущей способности при правильном подборе параметров текстуры.

Эффект объясняется возникновением микрогидродинамического давления на каждой ячейке текстуры: на входной по ходу вращения кромке углубления формируется локальный сходящийся клин, повышающий давление, а постоянное наличие жидкости в углублении обеспечивает резерв смазки в моменты пуска и останова и снижает абразивный износ за счёт удержания продуктов разрушения трибопары.

Количественная оценка перечисленных эффектов проводится численным интегрированием уравнения Рейнольдса~-- классического уравнения гидродинамики тонких смазочных слоёв. Решение этого уравнения для подшипника конечной длины позволяет получить двумерное поле давления и далее~-- ключевые интегральные показатели: подъёмную силу, коэффициент трения и объёмный расход смазки. В настоящей работе уравнение решается методом конечных разностей с итерационной схемой Гаусса--Зейделя и верхней релаксацией, что обеспечивает контролируемую сходимость на достаточно подробной сетке.

Практическая актуальность исследования определяется запросом отрасли на повышение надёжности и энергоэффективности насосно-компрессорного оборудования без принципиального изменения его конструкции~-- лазерное текстурирование вкладышей вписывается в существующие технологические маршруты и не требует переделки сопряжённых деталей.

Цель работы~-- получение и анализ трибологических характеристик радиального подшипника скольжения с дискретным микрорельефом методом численного моделирования.

В рамках поставленной цели решаются следующие задачи:
\begin{enumerate}
    \item рассмотреть теоретические основы гидродинамической смазки и обобщить опубликованные данные о влиянии регулярного микрорельефа на работу подшипников скольжения;
    \item реализовать численное решение уравнения Рейнольдса методом конечных разностей с итерационной процедурой Гаусса--Зейделя и оценкой сходимости;
    \item рассчитать поля давления и интегральные показатели как для гладкого, так и для текстурированного подшипника во всём диапазоне рабочих эксцентриситетов;
    \item сопоставить полученные результаты и количественно оценить эффект применения микрорельефа.
\end{enumerate}

В качестве объекта исследования выбран радиальный подшипник скольжения с \varDimpleTypeGenitive{} микрорельефом рабочей поверхности, соответствующий варианту~\varNumber{} (геометрия~\varGeomLetter).

Предмет исследования~-- закономерности изменения нагрузочной способности, коэффициента трения и расхода смазки при нанесении \varDimpleTypeGenitive{} углублений на поверхность вкладыша.

Использованы следующие методы: численное моделирование уравнения Рейнольдса методом конечных разностей, итерационная процедура Гаусса--Зейделя с верхней релаксацией, сравнительный анализ результирующих интегральных характеристик.

Работа имеет следующую структуру: введение, три главы, заключение, перечень использованных источников и приложение. В первой главе изложены теоретические основы гидродинамической смазки и приведён обзор работ по лазерному текстурированию поверхностей трения. Во второй главе сформулирована математическая модель и описана применённая численная схема. В третьей главе представлены результаты расчёта и проведён сравнительный анализ показателей гладкого и текстурированного подшипников.
